\section{Blood Viscosity from Red Blood Cell Partition Lag}
\label{sec:viscosity}

\subsection{Blood as a Suspension}

Blood consists of cellular elements (red blood cells, white blood cells, platelets) suspended in plasma. At physiological hematocrit (volume fraction of red cells) of 40-45\%, blood exhibits viscosity approximately 3-4 times that of water or plasma alone. This rheological behavior emerges from partition lag between red blood cells and surrounding plasma during flow.

From transport partition theory \cite{sachikonye2024transport}, all transport coefficients admit universal form
\begin{equation}
\Xi = \mathcal{N}^{-1} \sum_{i,j} \taulag_{ij} g_{ij}
\end{equation}

For viscosity, $\Xi = \mu$ and the normalization $\mathcal{N} = 1$. The partition lag $\taulag_{ij}$ represents time required for momentum transfer between particle $i$ and $j$, while coupling strength $g_{ij}$ quantifies the strength of their hydrodynamic interaction.

\subsection{Einstein's Relation for Dilute Suspensions}

For dilute suspensions where particles do not interact ($\phi \ll 1$, where $\phi$ is volume fraction), Einstein derived \cite{einstein1906new}
\begin{equation}
\mu = \mu_{0}(1 + 2.5\phi)
\label{eq:einstein}
\end{equation}
where $\mu_{0}$ is the viscosity of the suspending fluid.

This result emerges from considering the additional energy dissipation caused by flow disturbance around each particle. A single sphere of radius $a$ in shear flow $\dot{\gamma}$ dissipates power
\begin{equation}
\dot{E}_{\mathrm{sphere}} = 20\pi \mu_{0} a^{3} \dot{\gamma}^{2}
\end{equation}

Compared to the power dissipated by the fluid volume occupied by the sphere:
\begin{equation}
\dot{E}_{\mathrm{fluid}} = \frac{4}{3}\pi a^{3} \mu_{0} \dot{\gamma}^{2}
\end{equation}

The excess dissipation ratio is
\begin{equation}
\frac{\dot{E}_{\mathrm{sphere}}}{\dot{E}_{\mathrm{fluid}}} = \frac{20\pi \mu_{0} a^{3} \dot{\gamma}^{2}}{(4/3)\pi a^{3} \mu_{0} \dot{\gamma}^{2}} = 15
\end{equation}

But accounting for the volume excluded by the sphere (which cannot undergo deformation), the effective contribution to viscosity increase is factor 2.5 per unit volume fraction, yielding Equation \ref{eq:einstein}.

From the partition perspective, each suspended particle creates a momentum partition lag—fluid elements must transfer momentum around the obstacle rather than through direct shear. The partition lag scales with particle volume and number density.

\subsection{Concentrated Suspensions and the Krieger-Dougherty Relation}

At physiological hematocrit (∼45\%), red blood cell suspensions are far from dilute. Hydrodynamic interactions between cells become significant. The Krieger-Dougherty equation extends Einstein's relation to concentrated suspensions \cite{krieger1959mechanism}:
\begin{equation}
\mu = \mu_{0}\left(1 - \frac{\phi}{\phi_{\max}}\right)^{-[\mu]\phi_{\max}}
\end{equation}
where $\phi_{\max}$ is maximum packing fraction (approximately 0.64 for random close packing of spheres) and $[\mu] = 2.5$ is the intrinsic viscosity.

For red blood cells, which are deformable biconcave disks rather than rigid spheres, $\phi_{\max} \approx 0.95$ and $[\mu] \approx 2.5$ remains applicable. At $\phi = 0.45$:
\begin{equation}
\mu = \mu_{0}\left(1 - \frac{0.45}{0.95}\right)^{-2.5 \times 0.95} \approx \mu_{0}(0.526)^{-2.375} \approx 3.7\mu_{0}
\end{equation}

With plasma viscosity $\mu_{0} \approx 1.0-1.2$ cP, this predicts blood viscosity $\mu \approx 3.7-4.4$ cP, matching experimental measurements of 3-4 cP \cite{chien1970blood}.

\subsection{Red Blood Cell Deformability and Partition Lag}

Red blood cells are highly deformable, with membrane shear modulus approximately $\mu_{\mathrm{membrane}} \approx 6 \times 10^{-6}$ N/m. Under flow, cells elongate and align with streamlines, reducing resistance. The characteristic deformation time is
\begin{equation}
\taulag_{\mathrm{deform}} = \frac{\mu_{0} a}{\mu_{\mathrm{membrane}}} \approx \frac{(1.2 \times 10^{-3}~\mathrm{Pa \cdot s})(4 \times 10^{-6}~\mathrm{m})}{6 \times 10^{-6}~\mathrm{N/m}} \approx 8 \times 10^{-4}~\mathrm{s}
\end{equation}

For physiological shear rates $\dot{\gamma} \sim 100-1000$ s$^{-1}$ in larger vessels, $\taulag_{\mathrm{deform}} \cdot \dot{\gamma} \sim 0.1-1$, indicating cells experience significant deformation. This reduces effective viscosity compared to rigid particles.

If red blood cells were rigid disks, viscosity would be approximately 50\% higher. The deformability reduces partition lag by enabling cells to squeeze through constrictions and minimize flow disruption. Pathological conditions that reduce deformability (sickle cell disease, spherocytosis) correspondingly increase blood viscosity and vascular resistance.

\subsection{Fahraeus-Lindqvist Effect in Small Vessels}

In vessels with diameter less than approximately 300 μm, blood viscosity decreases with decreasing diameter—the Fahraeus-Lindqvist effect \cite{fahraeus1931viscosity}. This counterintuitive behavior reflects red cell migration toward the vessel center.

In Poiseuille flow, suspended particles experience lift force toward regions of lower shear rate. For red blood cells, this directs them away from walls toward the vessel centerline. The result is a cell-free layer near the wall, reducing effective viscosity.

From dimensional reduction in categorical fluid dynamics \cite{sachikonye2024fluid}, flow in a cylindrical vessel reduces to cross-sectional state plus axial transformation. The cell-free layer effectively increases the low-viscosity region near the wall, where most fluid passes.

The cell-free layer thickness $\delta$ scales with cell size:
\begin{equation}
\delta \approx 2-4~\mu\mathrm{m}
\end{equation}

For vessels with radius $r$, the effective viscosity is
\begin{equation}
\mu_{\mathrm{eff}} = \mu_{\mathrm{bulk}}\left(1 - \frac{\delta}{r}\right)^{4}
\end{equation}

The quartic dependence emerges from Poiseuille's $r^{4}$ scaling—flow rate is highly sensitive to radius changes.

For capillaries ($r \approx 3.5$ μm), $\delta/r \approx 0.6$:
\begin{equation}
\mu_{\mathrm{eff}} \approx \mu_{\mathrm{bulk}}(0.4)^{4} \approx 0.026\mu_{\mathrm{bulk}}
\end{equation}

The dramatic viscosity reduction in capillaries—by factor of approximately 40—substantially reduces microcirculatory resistance. Without this effect, capillary perfusion would require impractically high driving pressures.

\subsection{Plasma Composition and Baseline Viscosity}

Plasma itself is a non-Newtonian fluid due to dissolved proteins. Major constituents include:
\begin{itemize}
\item Albumin: 35-50 g/L
\item Globulins: 20-30 g/L
\item Fibrinogen: 2-4 g/L
\end{itemize}

These proteins increase viscosity through hydrodynamic volume and intermolecular interactions. The empirical relation for protein solutions is
\begin{equation}
\mu_{\mathrm{plasma}} = \mu_{\mathrm{water}} \exp(\alpha c)
\end{equation}
where $c$ is total protein concentration in g/dL and $\alpha \approx 0.015$ dL/g.

At physiological protein concentration $c \approx 7$ g/dL:
\begin{equation}
\mu_{\mathrm{plasma}} = (1.0~\mathrm{cP}) \exp(0.015 \times 7) \approx 1.1~\mathrm{cP}
\end{equation}

This establishes the baseline from which red blood cell contributions multiply.

Pathological increases in protein concentration (e.g., multiple myeloma with abnormal immunoglobulin production) elevate plasma viscosity to 2-3 cP, with corresponding increases in blood viscosity to 6-8 cP. This hyperviscosity syndrome impairs perfusion and can precipitate vascular occlusion.

\subsection{Temperature Dependence}

Blood viscosity depends strongly on temperature through the Arrhenius relation:
\begin{equation}
\mu(T) = \mu_{37} \exp\left(\frac{E_{a}}{\kB}\left(\frac{1}{T} - \frac{1}{310~\mathrm{K}}\right)\right)
\end{equation}

The activation energy $E_{a} \approx 0.1$ eV yields temperature coefficient approximately 2\% per degree Celsius. At hypothermia (30°C), viscosity increases approximately 30\%. At fever (40°C), viscosity decreases approximately 10\%.

This temperature sensitivity reflects the thermal dependence of molecular partition lag—higher temperature provides more kinetic energy for momentum transfer, reducing $\taulag$ and viscosity.

\subsection{Shear Rate Dependence and Non-Newtonian Behavior}

Blood exhibits shear-thinning behavior—viscosity decreases with increasing shear rate. At low shear ($\dot{\gamma} < 10$ s$^{-1}$), red cells aggregate into rouleaux (stacks), effectively increasing particle size and viscosity to 10-20 cP. At high shear ($\dot{\gamma} > 100$ s$^{-1}$), cells align and deform, reducing viscosity to 3-4 cP.

The Casson equation describes this behavior:
\begin{equation}
\sqrt{\tau} = \sqrt{\tau_{0}} + \sqrt{\mu_{\infty} \dot{\gamma}}
\end{equation}
where $\tau$ is shear stress, $\tau_{0} \approx 0.04$ dyn/cm$^{2}$ is yield stress (minimum stress to overcome rouleaux), and $\mu_{\infty} = 3.5$ cP is high-shear limiting viscosity.

This non-Newtonian character means blood viscosity is not a single value but depends on flow conditions. In large arteries with high shear rates, effective viscosity is approximately 3.5 cP. In veins with lower shear, it increases to 4-5 cP. In stagnant regions, it can exceed 10 cP.

\subsection{Experimental Validation}

Table \ref{tab:viscosity_validation} compares theoretical predictions with measurements across different conditions.

\begin{table}[h]
\centering
\caption{Blood viscosity: predicted versus measured values}
\label{tab:viscosity_validation}
\begin{tabular}{lcccc}
\toprule
\textbf{Condition} & \textbf{Hematocrit (\%)} & \textbf{Predicted (cP)} & \textbf{Measured (cP)} & \textbf{Reference} \\
\midrule
Normal male & 45 & 4.0 & $3.8-4.2$ & \cite{chien1970blood} \\
Normal female & 40 & 3.5 & $3.3-3.7$ & \cite{chien1970blood} \\
Anemia & 30 & 2.5 & $2.3-2.7$ & \cite{wells1962viscosity} \\
Polycythemia & 60 & 7.5 & $7.0-8.0$ & \cite{wells1962viscosity} \\
Capillary (est.) & 45 & 0.15 & $0.10-0.20$ & \cite{pries1994resistance} \\
\bottomrule
\end{tabular}
\end{table}

The close agreement validates the partition lag framework. Blood viscosity emerges from first principles as the collective partition lag of red blood cell-plasma interactions, modulated by cell deformability, aggregation, and the Fahraeus-Lindqvist effect in microvessels.
